\documentclass[11pt,letterpaper]{article}
\usepackage[T1]{fontenc}
\usepackage{lmodern}
\usepackage[margin=1.05in]{geometry}
\usepackage{amsmath,amssymb,amsthm,booktabs,longtable,microtype,xurl}
\usepackage[hidelinks,pdftitle={Conditional ZFC + I3 independence for a 380-state Turing machine},
  pdfauthor={James Chang}]{hyperref}

\newtheorem{theorem}{Theorem}[section]
\newtheorem{lemma}[theorem]{Lemma}
\newcommand{\ZFC}{\mathrm{ZFC}}
\newcommand{\PA}{\mathrm{PA}}
\newcommand{\Ithree}{\mathrm{I}_3}
\newcommand{\Con}{\operatorname{Con}}
\newcommand{\Halt}{\operatorname{Halt}}
\newcommand{\code}[1]{\texttt{\nolinkurl{#1}}}
\newcommand{\A}{\mathcal A}
\newcommand{\G}{\mathcal G}
\newcommand{\D}{\mathsf D}
\newcommand{\J}{\mathsf J}
\newcommand{\E}{\mathsf E}
\newcommand{\U}{\mathsf U}
\newcommand{\I}{\mathsf I}
\newcommand{\T}{\mathsf T}
\newcommand{\Sstage}{\mathsf S}
\newcommand{\ch}{\operatorname{ch}}
\newcommand{\loc}{\operatorname{h}}
\newcommand{\node}{\operatorname{node}}
\setlength{\emergencystretch}{2em}

\title{Conditional \texorpdfstring{$\ZFC+\Ithree$}{ZFC + I3} independence\texorpdfstring{\\}{ }for a 380-state Turing machine}
\author{James Chang}
\date{}

\begin{document}
\maketitle

\begin{abstract}
We formalize the conditional independence of blank-tape halting for a
literal two-symbol Turing machine with 380 working states. The machine
searches for a ground equational derivation of a distinguished equality.
Its nonhalting is equivalent to the syntactic consistency of
$\ZFC+\Ithree$. Under consistency, $\ZFC+\Ithree$ cannot prove the
machine's nonhalting. Under arithmetic $\Sigma_1$-soundness, it can
prove neither halting nor nonhalting. The unprovability of nonhalting
uses an internal consistency implication for a recursive arithmetic
theory whose theorems include the arithmetic consequences of
$\ZFC+\Ithree$. Lean verifies the theorem for the transition table
printed in the appendix.
\end{abstract}

\section{Introduction}\label{sec:result}

Let $M$ denote the binary Turing machine in Appendix~\ref{app:table},
and write $T=\ZFC+\Ithree$. We use the first-order formulation of
$\Ithree$ asserting that some $V_\lambda$, with $\lambda$ a nonzero
limit ordinal, admits a nonidentity elementary self-embedding in the
language of membership \cite{corazza}. The axiom uses a uniform
satisfaction formula for the set structure $V_\lambda$. In a
nonstandard model, the satisfaction assertion quantifies over the
formula codes of that model. The proof calculus is ordinary
first-order logic with equality, full Separation and Replacement, and
Choice.

Let $H_M$ be the closed set-theoretic sentence asserting the existence
of a finite computation of $M$ from the blank configuration to its
halt state. The finite history is coded by a function on an initial
segment of the natural numbers. Let $\Halt(M)$ denote actual halting
from the blank tape, with the number of steps taken in the standard
natural numbers. For arithmetic $\Sigma_1$-soundness, the formal statement
quantifies over a closed arithmetic $\Sigma_1$ formula $\theta$ and a
closed set-theoretic sentence $\psi$ that represents it. Representation
means that in every $T$-model, $\psi$ is true exactly when $\theta$ holds
on the natural numbers of the model. Soundness requires that if
$T\vdash\psi$ then $\mathbb N\models\theta$.

\begin{theorem}\label{thm:main}
For the machine $M$ the following hold.
\begin{enumerate}
\item $\neg\Halt(M)$ if and only if $\Con(T)$.
\item If $T$ is consistent, then $T\nvdash\neg H_M$.
\item If $T$ is arithmetically $\Sigma_1$-sound, then
      $T\nvdash H_M$ and $T\nvdash\neg H_M$.
\end{enumerate}
In every model of $T$, the sentence $H_M$ holds exactly when there is a
halting computation in the sense of that model.
\end{theorem}

The two steps in the proof of
Theorem~\ref{thm:main} concern different notions of computation.
The external equivalence uses standard finite rewrite sequences.
The unprovability result requires an implication that holds in
models of $T$ and therefore also applies to computations and proofs
whose lengths are nonstandard. Sections~\ref{sec:machine} and
\ref{sec:models} give the external reduction. Section~\ref{sec:internal}
constructs the internal implication and applies the second
incompleteness theorem.

\section{The literal machine and the ground search}\label{sec:machine}

The tape is doubly infinite, initially filled with zeros, and has
alphabet $\{0,1\}$. State $0$ is the start state, $0$ through $379$ are
working states, and $380$ is the separate halt label. A transition
$(q,b,w,d,q')$ reads $b$ in state $q$, writes $w$, moves one cell in
direction $d$, and enters $q'$. The only transition entering the halt
label is $(260,1,1,\mathrm R,380)$. The exact 760 rows are kept in
\code{data/zfc-i3-native-380.tm}, with SHA-256 digest
\begin{quote}\small
\nolinkurl{b9479f0a460656f93a13d0328a52dfdf410eaae7f1a8861f0b95141a0612c1cf}.
\end{quote}

The tape program implements a twelve-register program for one-step
equational rewrites.
A trial gives a proposed replacement word, a context position, the
orientation of a selected equation, and a dictionary assigning a
closed term to each of the equation's three available variables.
A trie reads the selectors for the equation. The verifier checks the
common prefix and suffix and the two instantiated sides of the
equation. An accepted trial replaces the current word by the proposal.
The machine halts when an accepted proposal is the code of $\J$.

To describe the word coding, put
$P(a,b)=\binom{a+b+1}{2}+b$. In the prefix token stream for a term,
a token of arity $a$ and tag $b$, followed by a stream with code $r$,
has code $P(a+1,P(b,r))$. The empty stream has code $0$.
The one-token words $\D$ and $\J$ have arity zero and tags one and
zero, so their word codes are $4$ and $1$, respectively.

The scheduler enumerates numeric trial streams. When a stream is
exhausted it starts a new one with incumbent $\D$. In the middle of a
stream the remaining code decreases. For a prescribed input $n$, the
termination argument for reaching that input uses the measure
\[
   (n-i)(n+1)+p,
\]
where $i\le n$ is the next input index and $p$ is the pending code.
Reading the next trial decreases $p$. Starting another stream
increases $i$ and loads a pending code bounded by the previous index.
This gives a decreasing measure until the requested input is reached,
unless a previous trial has already halted. Any finite valid rewrite
chain can be encoded consecutively in such a stream. As each trial is
read, the verifier consumes the corresponding fields and leaves the
continuation for the next trial.

\begin{lemma}[Exact search]\label{lem:search}
Let $\mathcal E_{19}$ be the equations defined in
Section~\ref{sec:equations}. Then
\[
  \Halt(M)\quad\Longleftrightarrow\quad
  \mathcal E_{19}\vdash_{\mathrm{eq}}\D=\J,
\]
where the derivations use ground substitution instances and contextual
replacement in either direction.
\end{lemma}

\begin{proof}
The compiler verification connects the literal configuration at each
control boundary to the register program that executes a trial. It
checks the physical decoder against a reference control tree and
checks the reference register program against the equation trie. The
literal blank run reaches its loader exit in 39,753 steps, and the
startup simulation connects that exit to the initial search store.
The equality of equation theories for the concrete trie rows and
$\mathcal E_{19}$ is proved by explicit equational derivations.

For soundness, the incumbent remains equationally equal to $\D$.
An unsuccessful trial contributes no new equation, and an accepted
trial checks one contextual rewrite of a ground substitution instance.
At a halt the accepted proposal is $\J$. For completeness, a derivation
of $\D=\J$ yields a finite chain of contextual rewrites. Encode that
chain as above and apply the scheduler's reachability argument. These
two directions use the checked control correspondence all the way
back to the literal blank tape.
\end{proof}

\section{The equational theory}\label{sec:equations}

Use constants $\D,\J,\U,\E,\I,\T,\Sstage$ and operations for Boolean
complement $\bar x$, converse $x^\smile$, relational composition
$x;y$, choice $\ch(x)$, a unary operation $\loc(x)$ and a ternary
constructor $\node(x,y,z)$. The machine uses NAND as a primitive,
and the term $x\wedge y$ abbreviates the complement of $x\mathbin{\mathrm{NAND}}y$.
In the relational interpretation $(x;y)(u,w)$ means
$\exists v\,(x(u,v)\land y(v,w))$.

The expanded presentation $\mathcal E_{21}$ has the two Meredith
equations \cite{meredith,robbins}
\begin{align*}
 \overline{(\bar a\wedge b)}\wedge a &= a,\\
 \overline{(\bar a\wedge b)}\wedge(c\wedge b)
   &= b\wedge(c\wedge a),
\end{align*}
and the two defining equations
\[
 \T=\overline{\E^\smile;\bar\E},
 \qquad
 \Sstage=\ch(\T)^\smile;\ch(\T).
\]
Here $\T$ represents inclusion, and $\Sstage$ represents the diagonal
on hierarchy stages in the model built from a rank embedding.
The remaining equations are as follows. Consecutive products are
associated to the left where parentheses are omitted.
\begin{align*}
 a;(b;c)&=(a;b);c, &
 a;\I&=a,\\
 (a^\smile)^\smile&=a, &
 (a;b)\wedge(a;(c\wedge b))&=a;(c\wedge b),\\
 (a\wedge b)^\smile&=a^\smile\wedge b^\smile, &
 (a^\smile;\overline{a;b})\wedge b&=\bar\U,\\
 \T\wedge\T^\smile&=\I, &
 \ch(a)\wedge a&=\ch(a),\\
 (\ch(a)^\smile;\ch(a))\wedge\bar\I&=\bar\U, &
 \ch(a);\U&=a;\U,\\
 \D&=\I\wedge(\J;\U), &
 \loc(\I)&=\E;\D.
\end{align*}
Writing $r=\E^\smile;\ch(a)$ as a term abbreviation, the other five
equations are
\begin{gather*}
 (\overline{r;\bar\E}\wedge\overline{\bar r;\E});\U=\U,\\
 \Sstage;\E^\smile=\Sstage;\E^\smile;\Sstage;\T^\smile,\\
 \D=\I\wedge(\E;(\Sstage\wedge\D);\U),\\
 \loc(\node(a,b,c))
   =(\loc(a);\loc(b)^\smile)
        \wedge(\loc(b)\mathbin{\mathrm{NAND}}\loc(c)),\\
 (\J^\smile;\loc(a);\J)\wedge\overline{\loc(a)}=\bar\U.
\end{gather*}
The equations are universally closed in the variables displayed.

To obtain $\mathcal E_{19}$, substitute the defining terms for $\T$
and $\Sstage$, remove the definition equations, and cancel common
leading complements on both sides of a remaining equation. The
Meredith equations stay in the presentation and give complement
involution. Thus each cancellation is reversible. Every model of
$\mathcal E_{19}$ admits the corresponding expansion to a model of
$\mathcal E_{21}$, and restriction gives the converse. In particular,
the elimination preserves ground derivability of $\D=\J$. The
physical trie uses equation instances from a fixed presentation with
reassociated products, side reversals, and variable renamings.
The Lean equivalence for those rows also checks the few replacements
by equationally equivalent rows.

\section{Equational underivability and consistency}\label{sec:models}

\begin{lemma}\label{lem:equational}
$\mathcal E_{19}\nvdash_{\mathrm{eq}}\D=\J$ if and only if $T$ is
syntactically consistent.
\end{lemma}

\begin{proof}
First suppose that $\D=\J$ is underivable. Quotient the ground term
algebra by equational derivability and use the expansion to
$\mathcal E_{21}$. The two constants remain distinct in this quotient.
In the Boolean algebra of subidentities below $\I$, choose an
ultrafilter containing the domain support of a nonzero separating
piece of $\D$ and $\J$. More explicitly, a separating piece is one of
$\D\wedge\bar\J$ and $\J\wedge\bar\D$, and the support of $r$ is
$\I\wedge(r;\U)$. A countable enumeration of the ground classes
allows the filter to be constructed by successive nonzero meets.
The representation takes equivalence classes of local functional
elements as points, where agreement on a support in the ultrafilter
identifies two local functions. Relational application to those
points gives a square in which the images of $\D$ and $\J$ differ.

Restrict to the points that have singleton names in the representation.
The equations for choice, image and the stage operations give the
needed closure on these points. The pairing construction supplies
projections for translating arbitrary first-order formulas with
parameters. The resulting interpretation of $\E$ satisfies
Extensionality, Empty Set, Pairing, Union, Power Set and Choice.
The definable relation and image constructions give the full
Separation and Replacement schemas.

The interpretation of $\D$ is an active subidentity, and the
interpretation of $\J$ is a function on its active points. Its
difference from $\D$ gives a moved point. On the active domain, the
local constructor equations allow membership formulas to be expressed
by the relations $\loc(a)$. The preservation equation for $\J$ then
gives local elementarity. The stage law makes the rank analysis
possible before all the set axioms have been recovered. In particular,
the rank argument supplies Foundation, and the moved-point argument
supplies an active inductive set and an active $\omega$.

Choose a critical point and form its critical sequence by recursion
on that internal $\omega$. If $\kappa_n$ denotes the resulting
sequence of heights, the limit is $\lambda=\sup_{n\in\omega}\kappa_n$.
The corresponding increasing stages have union $V_\lambda$.
Replacement constructs the graph of the restriction of $\J$ to this
union, and the next-stage property ensures that its range is in the
union. A coded assignment of parameters from the union is bounded
in one earlier stage. Satisfaction on the union is obtained by set
recursion on formula codes. Transport of satisfaction across a
successive pair of stages gives preservation on the union. The
critical point is moved, so the restriction witnesses the literal
$\Ithree$ axiom. This produces a model of $T$, and soundness implies
$\Con(T)$.

Conversely, use completeness to obtain a countable model of $T$.
The construction expands it by a generic global order such that
expanded formulas have Separation and Replacement and definable
nonempty successor classes have least elements. The membership
relation remains the one in the original model. Take the algebra of
definable binary relations with parameters in this expanded
structure. Interpret $\E$ as membership, $\I$ as the diagonal and
$\U$ as the universal relation. Interpret $\D$ as the diagonal on the
rank of the given embedding and $\J$ as the graph of that embedding.
For choice use least eligible successors, with a stage selector on
the inclusion relation. Interpret the local operations by evaluation
of finite ternary trees on the rank, starting from membership
restricted to that rank. Elementarity gives the local preservation
equation. These interpretations satisfy $\mathcal E_{21}$ and
separate $\D$ from $\J$. Equational soundness gives the desired
underivability.
\end{proof}

Combining Lemmas~\ref{lem:search} and~\ref{lem:equational} gives the
first assertion of Theorem~\ref{thm:main}.

\section{The internal consistency implication}\label{sec:internal}

For an arithmetic sentence $\varphi$, write $\varphi^{\mathrm{set}}$
for its reading on the natural numbers of a set model. We construct
a recursive arithmetic theory $\A$ with the properties
\begin{enumerate}
\item $\A$ extends $\mathrm I\Sigma_1$, and $\Con(T)$ implies $\Con(\A)$.
\item If $T\vdash\varphi^{\mathrm{set}}$, then $\A\vdash\varphi$.
\item $T\vdash\neg H_M\to\Con(\A)^{\mathrm{set}}$.
\end{enumerate}
Here $\Con(\A)$ is the consistency sentence for the standard
arithmetized provability predicate of the specified recursive
presentation.

\subsection{A finite theory of separating algebras}

Expand the language of $\mathcal E_{21}$ by a predicate for a
subidentity ultrafilter. Add the ultrafilter axioms and the assertion
that it contains the support of a separating piece of $\D$ and $\J$.
Convert this theory to a finite relational theory $\G$ by replacing
function symbols by their graphs and adding totality and uniqueness.
The relational form allows term and assignment operations in the
coded proof calculus to use variable terms.

A model of $\G$ gives a named square by the representation in the
first half of Lemma~\ref{lem:equational}. Reading the natural numbers
in the resulting set model interprets arithmetic in $\G$. Denote the
translation of an arithmetic sentence by $\varphi^\dagger$. Set-theoretic
soundness and first-order completeness give
\[
  T\vdash\varphi^{\mathrm{set}}
  \quad\Longrightarrow\quad
  \G\vdash\varphi^\dagger .
\]
The proof of this implication checks the translation against the
operations on the natural numbers of the represented model.

\subsection{A recursive arithmetic presentation}

Use the sentences $\varphi$ for which $\G\vdash\varphi^\dagger$.
The recursive presentation pads $\varphi$ by a conjunction of
tautologies indexed by a code $p$ of a $\G$-proof of
$\varphi^\dagger$. In the padded sentence, the candidate witness is
bounded by its code, so axiom recognition checks a bounded witness
for the translation and proof relation. The resulting presentation has
a $\Delta_1$ axiom predicate in the coding used for incompleteness.
A padded axiom proves its unpadded conjunct. In particular, $\A$
absorbs arithmetic consequences of $T$. Since $T$ proves the readings
of PA theorems on its natural numbers, $\A$ extends PA. If $T$ is
consistent, the separating term algebra gives a model of $\G$.
Its interpreted natural numbers satisfy every admitted padded axiom,
which gives consistency of $\A$.

\subsection{Internal nonhalting and coded proof soundness}

Fix any model $N$ of $T$ in which $H_M$ is false. Work with its own
natural numbers. Codes for ground terms and finite rewrite sequences
are arithmetic in this model. The internal replay theorem proves
that collapse of the two ground classes $\D$ and $\J$ would give a
halting computation in the model. The proof uses arithmetic induction
to replay a coded sequence and to reach the corresponding scheduler
input. Its induction applies to every internally finite code.
The two classes are therefore distinct.

Take least representatives of the ground classes and construct the
separating filter by a coded descending sequence. The resulting
domain and relation tables are sets of $N$. They give a set-coded
model of $\G$ in $N$. Set recursion in $N$ supplies a truth table
on its arithmetic formula codes for this structure. Induction on
coded sequent derivations proves soundness with respect to that
table. In particular, there is no coded $\G$-refutation in $N$.
The verification treats the truth table as a set parameter in the
induction argument.

The translation of arithmetic proofs to the relational language is
also checked on internal codes. An admitted $\A$ axiom comes with a
$\G$-proof of its translated arithmetic conjunct. The checked
translation of derivation rules therefore takes an $\A$-refutation
in $N$ to a $\G$-refutation in $N$. This proves
$\Con(\A)^{\mathrm{set}}$ in every model of $T$ satisfying
$\neg H_M$. Completeness for the source calculus proves the
internal implication.

\begin{proof}[Proof of the remaining assertions of Theorem~\ref{thm:main}]
Suppose $T$ is consistent and $T\vdash\neg H_M$. The internal
implication gives $T\vdash\Con(\A)^{\mathrm{set}}$. Absorption yields
$\A\vdash\Con(\A)$. Since $\A$ is a consistent recursive extension of
$\mathrm I\Sigma_1$, its standard provability predicate satisfies the
derivability conditions used in G\"odel's second incompleteness
theorem \cite{goedel}. This is a contradiction.

The finite run predicate of the literal machine has a closed
arithmetic $\Sigma_1$ representation. In $\mathbb N$ it expresses
actual halting, and its reading on the natural numbers in each
$T$-model agrees with $H_M$. If $T$ is arithmetically
$\Sigma_1$-sound, it is consistent, so the external equivalence gives
actual nonhalting. A $T$-proof of $H_M$ would contradict soundness.
The preceding argument excludes a proof of $\neg H_M$.
\end{proof}

\section{Formal verification}\label{sec:lean}

The Lean project contains the exact table, a standalone Challenge
statement and a Solution. The Solution uses the closed sentence
\code{PalomarI3380.Logic.haltingSentence}. The principal comparison,
\code{PalomarI3380.exactClaim}, includes the representation of halting
in each model and the two conditional conclusions. The external
equivalence with consistency is the separate Solution theorem
\code{PalomarI3380.literal_nonhalting_iff_consistency}.

The link from the tape machine to the ground calculus is
\code{Compiler.Source.exact_machine_halts_iff_derives}.
The model-valued converse used in the internal argument is
\code{Calibration.TermWords.model_halts_of_distinction_failure}.
Soundness of coded graph proofs and the internal consistency
implication are verified in \code{Calibration/InternalGraphSoundness.lean}
and \code{Calibration/InternalConsistency.lean}.

The project uses Foundation \cite{foundation} for arithmetization
and incompleteness. The build and axiom audit passed using Lean
4.37.0-rc1. Comparator accepted the full principal declaration in
both Lean and the bundled NanoDa checker in a local macOS run.
The declaration depends on \code{propext}, \code{Classical.choice}
and \code{Quot.sound}. The reproduction instructions give the
exact commands. The paper and Lean development were
prepared with assistance from Codex.

\begingroup
\small
\begin{thebibliography}{9}
\setlength{\itemsep}{0pt}
\bibitem{corazza}
P. Corazza,
The gap between $\mathrm I_3$ and the wholeness axiom,
\emph{Fundamenta Mathematicae} \textbf{179} (2003), 43--60.
\href{https://doi.org/10.4064/fm179-1-4}{doi:10.4064/fm179-1-4}.

\bibitem{meredith}
C. A. Meredith and A. N. Prior,
Equational logic,
\emph{Notre Dame Journal of Formal Logic} \textbf{9} (1968), 212--226.
\href{https://doi.org/10.1305/ndjfl/1093893457}{doi:10.1305/ndjfl/1093893457}.

\bibitem{robbins}
W. Truszkowska and A. Grabowski,
On the two short axiomatizations of ortholattices,
\emph{Journal of Formalized Mathematics} \textbf{15} (2003),
Mizar article ROBBINS2.
\url{https://mizar.uwb.edu.pl/JFM/Vol15/robbins2.html}.

\bibitem{goedel}
K. G\"odel,
\"Uber formal unentscheidbare S\"atze der Principia Mathematica
und verwandter Systeme I,
\emph{Monatshefte f\"ur Mathematik und Physik} \textbf{38} (1931),
173--198.
\href{https://doi.org/10.1007/BF01700692}{doi:10.1007/BF01700692}.

\bibitem{foundation}
S. Saitou and M. Noguchi,
\emph{Formalized Formal Logic}, Lean library.
\url{https://github.com/FormalizedFormalLogic/Foundation}.
\end{thebibliography}
\endgroup

\clearpage
\appendix
\section{The transition table}\label{app:table}

In each group the two entries are for reading $0$ and $1$.
An entry of the form \texttt{1R023} writes $1$, moves right and
enters state $23$. State $0$ starts the blank run and $380$ halts it.

\begingroup
\small
\setlength{\tabcolsep}{5pt}
\begin{longtable}{@{}rll@{\qquad}rll@{\qquad}rll@{\qquad}rll@{}}
\toprule
$q$ & $0$ & $1$ & $q$ & $0$ & $1$ & $q$ & $0$ & $1$ & $q$ & $0$ & $1$\\
\midrule
\endfirsthead
\toprule
$q$ & $0$ & $1$ & $q$ & $0$ & $1$ & $q$ & $0$ & $1$ & $q$ & $0$ & $1$\\
\midrule
\endhead
\bottomrule
\endfoot
0 & \texttt{1R001} & \texttt{0L002} & 95 & \texttt{1L062} & \texttt{1R080} & 190 & \texttt{0R176} & \texttt{1R183} & 285 & \texttt{0R175} & \texttt{1R231} \\
1 & \texttt{0R002} & \texttt{0R001} & 96 & \texttt{0R042} & \texttt{1R049} & 191 & \texttt{0R171} & \texttt{1R207} & 286 & \texttt{0R014} & \texttt{1R296} \\
2 & \texttt{1L003} & \texttt{0R004} & 97 & \texttt{0R042} & \texttt{1R051} & 192 & \texttt{0R156} & \texttt{1R173} & 287 & \texttt{0R297} & \texttt{1R298} \\
3 & \texttt{0L004} & \texttt{0L005} & 98 & \texttt{0R109} & \texttt{1R110} & 193 & \texttt{0R081} & \texttt{1R208} & 288 & \texttt{0R154} & \texttt{1R249} \\
4 & \texttt{1L000} & \texttt{0L001} & 99 & \texttt{0R111} & \texttt{0L047} & 194 & \texttt{0R209} & \texttt{1R210} & 289 & \texttt{0R299} & \texttt{1R300} \\
5 & \texttt{0L006} & \texttt{1R007} & 100 & \texttt{0R060} & \texttt{1R028} & 195 & \texttt{0R057} & \texttt{1R165} & 290 & \texttt{0R081} & \texttt{1R301} \\
6 & \texttt{0L005} & \texttt{0L005} & 101 & \texttt{0R096} & \texttt{0L047} & 196 & \texttt{0R084} & \texttt{1R211} & 291 & \texttt{0R302} & \texttt{1R303} \\
7 & \texttt{0R008} & \texttt{1R009} & 102 & \texttt{0R112} & \texttt{0L047} & 197 & \texttt{0R022} & \texttt{1R212} & 292 & \texttt{0R081} & \texttt{1R304} \\
8 & \texttt{0R010} & \texttt{1R011} & 103 & \texttt{0R113} & \texttt{1R114} & 198 & \texttt{0R213} & \texttt{1R145} & 293 & \texttt{0R131} & \texttt{1R187} \\
9 & \texttt{0R012} & \texttt{1R013} & 104 & \texttt{0L115} & \texttt{1L038} & 199 & \texttt{0R214} & \texttt{1R203} & 294 & \texttt{0R305} & \texttt{1R306} \\
10 & \texttt{0L012} & \texttt{1L014} & 105 & \texttt{0R116} & \texttt{0R090} & 200 & \texttt{0R215} & \texttt{1R142} & 295 & \texttt{0R307} & \texttt{1R308} \\
11 & \texttt{0R015} & \texttt{1R016} & 106 & \texttt{0L117} & \texttt{1L091} & 201 & \texttt{0R216} & \texttt{1R215} & 296 & \texttt{0R309} & \texttt{1R310} \\
12 & \texttt{0R017} & \texttt{1R018} & 107 & \texttt{0L042} & \texttt{1R116} & 202 & \texttt{0R217} & \texttt{1R123} & 297 & \texttt{0R081} & \texttt{1R311} \\
13 & \texttt{0R019} & \texttt{1R020} & 108 & \texttt{1L118} & \texttt{0L062} & 203 & \texttt{0R218} & \texttt{1R219} & 298 & \texttt{0R275} & \texttt{1R312} \\
14 & \texttt{0R021} & \texttt{1R022} & 109 & \texttt{0R068} & \texttt{1R119} & 204 & \texttt{0R209} & \texttt{1R215} & 299 & \texttt{0R081} & \texttt{1R313} \\
15 & \texttt{0R023} & \texttt{1R024} & 110 & \texttt{0R029} & \texttt{1R085} & 205 & \texttt{0R220} & \texttt{1R221} & 300 & \texttt{0R314} & \texttt{1R315} \\
16 & \texttt{0R025} & \texttt{1R026} & 111 & \texttt{0R042} & \texttt{1R120} & 206 & \texttt{0R222} & \texttt{1R171} & 301 & \texttt{0R316} & \texttt{0L036} \\
17 & \texttt{0R027} & \texttt{1R028} & 112 & \texttt{0R121} & \texttt{1R121} & 207 & \texttt{0R185} & \texttt{1R128} & 302 & \texttt{0R317} & \texttt{1R192} \\
18 & \texttt{0R029} & \texttt{1R030} & 113 & \texttt{0L122} & \texttt{1L123} & 208 & \texttt{0R223} & \texttt{0L036} & 303 & \texttt{0R317} & \texttt{1R318} \\
19 & \texttt{0L031} & \texttt{1R032} & 114 & \texttt{0R124} & \texttt{1R125} & 209 & \texttt{0R207} & \texttt{1R224} & 304 & \texttt{0R122} & \texttt{0L036} \\
20 & \texttt{0R033} & \texttt{1R034} & 115 & \texttt{0L126} & \texttt{1L038} & 210 & \texttt{0R072} & \texttt{1R225} & 305 & \texttt{0R081} & \texttt{1R319} \\
21 & \texttt{0R035} & \texttt{0L036} & 116 & \texttt{0L059} & \texttt{1R116} & 211 & \texttt{0R137} & \texttt{1R226} & 306 & \texttt{0R246} & \texttt{1R320} \\
22 & \texttt{0R037} & \texttt{0L038} & 117 & \texttt{0L117} & \texttt{0L036} & 212 & \texttt{0R035} & \texttt{0L038} & 307 & \texttt{0R321} & \texttt{1R322} \\
23 & \texttt{0R039} & \texttt{1R040} & 118 & \texttt{0R108} & \texttt{1L118} & 213 & \texttt{0R048} & \texttt{1R227} & 308 & \texttt{0R323} & \texttt{1R141} \\
24 & \texttt{0R029} & \texttt{1R041} & 119 & \texttt{0R127} & \texttt{0L047} & 214 & \texttt{0R175} & \texttt{1R228} & 309 & \texttt{0R035} & \texttt{1R324} \\
25 & \texttt{0R042} & \texttt{1R043} & 120 & \texttt{0R042} & \texttt{1R066} & 215 & \texttt{0R198} & \texttt{1R229} & 310 & \texttt{0R147} & \texttt{1R325} \\
26 & \texttt{0R044} & \texttt{1R045} & 121 & \texttt{0R098} & \texttt{1R128} & 216 & \texttt{0R154} & \texttt{1R230} & 311 & \texttt{0R326} & \texttt{0L036} \\
27 & \texttt{0R042} & \texttt{1R039} & 122 & \texttt{0R129} & \texttt{1R130} & 217 & \texttt{0R185} & \texttt{1R185} & 312 & \texttt{0R249} & \texttt{1R183} \\
28 & \texttt{0R046} & \texttt{0L047} & 123 & \texttt{0R131} & \texttt{1R132} & 218 & \texttt{0R175} & \texttt{1R010} & 313 & \texttt{0R327} & \texttt{0L036} \\
29 & \texttt{0R042} & \texttt{1R048} & 124 & \texttt{0R081} & \texttt{0L038} & 219 & \texttt{0R231} & \texttt{1R232} & 314 & \texttt{0L328} & \texttt{1L278} \\
30 & \texttt{0R049} & \texttt{0L047} & 125 & \texttt{0R012} & \texttt{1R133} & 220 & \texttt{0R183} & \texttt{1R204} & 315 & \texttt{0R154} & \texttt{1R329} \\
31 & \texttt{0R050} & \texttt{1R018} & 126 & \texttt{0R134} & \texttt{1L038} & 221 & \texttt{0R233} & \texttt{1R234} & 316 & \texttt{0R253} & \texttt{1R330} \\
32 & \texttt{0R051} & \texttt{0L047} & 127 & \texttt{0R066} & \texttt{1R135} & 222 & \texttt{0R187} & \texttt{1R142} & 317 & \texttt{0R154} & \texttt{1R314} \\
33 & \texttt{0R052} & \texttt{1R019} & 128 & \texttt{0R136} & \texttt{1R137} & 223 & \texttt{0R235} & \texttt{1R236} & 318 & \texttt{0R323} & \texttt{1R302} \\
34 & \texttt{0R053} & \texttt{1R054} & 129 & \texttt{0R081} & \texttt{1R138} & 224 & \texttt{0R170} & \texttt{1R128} & 319 & \texttt{0R331} & \texttt{0L036} \\
35 & \texttt{0R042} & \texttt{1R055} & 130 & \texttt{0R139} & \texttt{1R140} & 225 & \texttt{0R237} & \texttt{1R238} & 320 & \texttt{0R332} & \texttt{1R333} \\
36 & \texttt{0L056} & \texttt{0L036} & 131 & \texttt{0R141} & \texttt{1R142} & 226 & \texttt{0R239} & \texttt{1R240} & 321 & \texttt{0R207} & \texttt{1R334} \\
37 & \texttt{0R042} & \texttt{1R057} & 132 & \texttt{0R143} & \texttt{1R144} & 227 & \texttt{0R174} & \texttt{0L047} & 322 & \texttt{0R207} & \texttt{1R192} \\
38 & \texttt{0L058} & \texttt{1L038} & 133 & \texttt{0R145} & \texttt{1R146} & 228 & \texttt{0R212} & \texttt{1R231} & 323 & \texttt{0R307} & \texttt{1R142} \\
39 & \texttt{0R042} & \texttt{1R059} & 134 & \texttt{0R134} & \texttt{1R007} & 229 & \texttt{0R214} & \texttt{1R199} & 324 & \texttt{0R284} & \texttt{0L047} \\
40 & \texttt{0R060} & \texttt{0L047} & 135 & \texttt{0R021} & \texttt{1R147} & 230 & \texttt{0R154} & \texttt{1R241} & 325 & \texttt{0R335} & \texttt{1R128} \\
41 & \texttt{0R025} & \texttt{0L047} & 136 & \texttt{0R148} & \texttt{1R149} & 231 & \texttt{0L135} & \texttt{1R174} & 326 & \texttt{0R336} & \texttt{1R189} \\
42 & \texttt{1R061} & \texttt{0L062} & 137 & \texttt{0R150} & \texttt{1R151} & 232 & \texttt{0R242} & \texttt{1R243} & 327 & \texttt{0R337} & \texttt{1R338} \\
43 & \texttt{1R063} & \texttt{1R043} & 138 & \texttt{0R152} & \texttt{0L036} & 233 & \texttt{0R222} & \texttt{1R244} & 328 & \texttt{0R144} & \texttt{1R339} \\
44 & \texttt{0R064} & \texttt{0L038} & 139 & \texttt{0R153} & \texttt{1R142} & 234 & \texttt{0R204} & \texttt{1R183} & 329 & \texttt{0R340} & \texttt{1R341} \\
45 & \texttt{0R065} & \texttt{1R065} & 140 & \texttt{0R154} & \texttt{1R155} & 235 & \texttt{0R081} & \texttt{1R245} & 330 & \texttt{0R154} & \texttt{1R342} \\
46 & \texttt{0R066} & \texttt{1R067} & 141 & \texttt{0R156} & \texttt{1R157} & 236 & \texttt{0R246} & \texttt{1R247} & 331 & \texttt{0R343} & \texttt{1R344} \\
47 & \texttt{0L047} & \texttt{0L038} & 142 & \texttt{0R158} & \texttt{1R159} & 237 & \texttt{0R081} & \texttt{1R248} & 332 & \texttt{0R081} & \texttt{1R345} \\
48 & \texttt{0R042} & \texttt{1R068} & 143 & \texttt{0L160} & \texttt{1L161} & 238 & \texttt{0R249} & \texttt{1R250} & 333 & \texttt{0R217} & \texttt{1R346} \\
49 & \texttt{0R042} & \texttt{1R025} & 144 & \texttt{0R156} & \texttt{1R156} & 239 & \texttt{0R048} & \texttt{0L038} & 334 & \texttt{0R288} & \texttt{1R192} \\
50 & \texttt{0R055} & \texttt{0L036} & 145 & \texttt{0L031} & \texttt{1R162} & 240 & \texttt{0R037} & \texttt{1R251} & 335 & \texttt{0R156} & \texttt{1R285} \\
51 & \texttt{0R042} & \texttt{1R069} & 146 & \texttt{0R163} & \texttt{1R164} & 241 & \texttt{0R187} & \texttt{1R183} & 336 & \texttt{0R081} & \texttt{1R347} \\
52 & \texttt{0R037} & \texttt{1R070} & 147 & \texttt{0R029} & \texttt{1R165} & 242 & \texttt{0R082} & \texttt{1R252} & 337 & \texttt{0R081} & \texttt{1R348} \\
53 & \texttt{0R064} & \texttt{1R071} & 148 & \texttt{0R057} & \texttt{1R028} & 243 & \texttt{0R074} & \texttt{0L038} & 338 & \texttt{0R187} & \texttt{1R349} \\
54 & \texttt{0R072} & \texttt{1R073} & 149 & \texttt{0L098} & \texttt{1R071} & 244 & \texttt{0R253} & \texttt{1R254} & 339 & \texttt{0R350} & \texttt{1R113} \\
55 & \texttt{0R042} & \texttt{1R074} & 150 & \texttt{0R035} & \texttt{1R166} & 245 & \texttt{0R255} & \texttt{0L036} & 340 & \texttt{0R217} & \texttt{1R183} \\
56 & \texttt{1L038} & \texttt{0L056} & 151 & \texttt{0R055} & \texttt{1R167} & 246 & \texttt{0R154} & \texttt{1R072} & 341 & \texttt{0R351} & \texttt{1R352} \\
57 & \texttt{0R042} & \texttt{1R064} & 152 & \texttt{0R154} & \texttt{1R168} & 247 & \texttt{0R256} & \texttt{1R257} & 342 & \texttt{0R321} & \texttt{1R353} \\
58 & \texttt{0L075} & \texttt{1L038} & 153 & \texttt{0R156} & \texttt{1R169} & 248 & \texttt{0R258} & \texttt{0L036} & 343 & \texttt{0R081} & \texttt{1R354} \\
59 & \texttt{1R076} & \texttt{1L077} & 154 & \texttt{0R170} & \texttt{1R156} & 249 & \texttt{0R185} & \texttt{1R170} & 344 & \texttt{0R141} & \texttt{1R355} \\
60 & \texttt{0R049} & \texttt{1R066} & 155 & \texttt{0R171} & \texttt{1R153} & 250 & \texttt{0R072} & \texttt{1R259} & 345 & \texttt{0R356} & \texttt{0L036} \\
61 & \texttt{0R061} & \texttt{1R078} & 156 & \texttt{0R172} & \texttt{1R128} & 251 & \texttt{0R260} & \texttt{1R261} & 346 & \texttt{0R357} & \texttt{1R358} \\
62 & \texttt{0L079} & \texttt{1L062} & 157 & \texttt{0R172} & \texttt{1R173} & 252 & \texttt{0R064} & \texttt{1R262} & 347 & \texttt{0R202} & \texttt{0L036} \\
63 & \texttt{0R063} & \texttt{1R080} & 158 & \texttt{0R057} & \texttt{1R174} & 253 & \texttt{0R209} & \texttt{1R207} & 348 & \texttt{0R359} & \texttt{0L036} \\
64 & \texttt{0R042} & \texttt{1R081} & 159 & \texttt{0R149} & \texttt{1R175} & 254 & \texttt{0L233} & \texttt{1L263} & 349 & \texttt{0R207} & \texttt{1R360} \\
65 & \texttt{0R082} & \texttt{1R082} & 160 & \texttt{0R176} & \texttt{1R177} & 255 & \texttt{0R264} & \texttt{1R265} & 350 & \texttt{0R288} & \texttt{1R361} \\
66 & \texttt{0R042} & \texttt{1R083} & 161 & \texttt{0R178} & \texttt{1R179} & 256 & \texttt{0R081} & \texttt{1R266} & 351 & \texttt{0R314} & \texttt{1R207} \\
67 & \texttt{0R050} & \texttt{1R084} & 162 & \texttt{0R083} & \texttt{0L047} & 257 & \texttt{0R187} & \texttt{1R267} & 352 & \texttt{0R217} & \texttt{1R143} \\
68 & \texttt{0R042} & \texttt{1R037} & 163 & \texttt{0R072} & \texttt{1R180} & 258 & \texttt{0R268} & \texttt{1R269} & 353 & \texttt{0R171} & \texttt{1R362} \\
69 & \texttt{0R042} & \texttt{1R082} & 164 & \texttt{0R124} & \texttt{1R181} & 259 & \texttt{0R270} & \texttt{1R271} & 354 & \texttt{0R205} & \texttt{0L036} \\
70 & \texttt{0R085} & \texttt{0L047} & 165 & \texttt{0R172} & \texttt{0L047} & 260 & \texttt{0R272} & \texttt{1R380} & 355 & \texttt{0R246} & \texttt{1R363} \\
71 & \texttt{0R069} & \texttt{0L047} & 166 & \texttt{0R182} & \texttt{0L047} & 261 & \texttt{0R065} & \texttt{1R273} & 356 & \texttt{0R241} & \texttt{1R160} \\
72 & \texttt{0R012} & \texttt{1R086} & 167 & \texttt{0R022} & \texttt{0L047} & 262 & \texttt{0R274} & \texttt{0L047} & 357 & \texttt{0R254} & \texttt{1R340} \\
73 & \texttt{0R087} & \texttt{1R088} & 168 & \texttt{0R183} & \texttt{1R184} & 263 & \texttt{0R275} & \texttt{1R276} & 358 & \texttt{0R224} & \texttt{1R364} \\
74 & \texttt{0R042} & \texttt{1R029} & 169 & \texttt{0R172} & \texttt{1R170} & 264 & \texttt{0R081} & \texttt{1R277} & 359 & \texttt{0R171} & \texttt{1R365} \\
75 & \texttt{0L089} & \texttt{1L038} & 170 & \texttt{0R172} & \texttt{1R185} & 265 & \texttt{0R278} & \texttt{1R279} & 360 & \texttt{0R230} & \texttt{1R143} \\
76 & \texttt{0R076} & \texttt{0R090} & 171 & \texttt{0R186} & \texttt{1R187} & 266 & \texttt{0R280} & \texttt{0L036} & 361 & \texttt{0R366} & \texttt{1R366} \\
77 & \texttt{0L091} & \texttt{1R092} & 172 & \texttt{0R042} & \texttt{1R111} & 267 & \texttt{0R072} & \texttt{1R281} & 362 & \texttt{0R192} & \texttt{1R367} \\
78 & \texttt{0R093} & \texttt{1R078} & 173 & \texttt{0R172} & \texttt{1R169} & 268 & \texttt{0R081} & \texttt{1R282} & 363 & \texttt{0R368} & \texttt{1R369} \\
79 & \texttt{0L094} & \texttt{1L062} & 174 & \texttt{0R120} & \texttt{1L100} & 269 & \texttt{0R171} & \texttt{1R206} & 364 & \texttt{0R340} & \texttt{1R370} \\
80 & \texttt{0R095} & \texttt{1R080} & 175 & \texttt{0R074} & \texttt{1L098} & 270 & \texttt{0R081} & \texttt{1R283} & 365 & \texttt{0R191} & \texttt{1R186} \\
81 & \texttt{0R042} & \texttt{1R027} & 176 & \texttt{0R154} & \texttt{1R188} & 271 & \texttt{0R254} & \texttt{1R263} & 366 & \texttt{0R186} & \texttt{1R307} \\
82 & \texttt{0R042} & \texttt{1R096} & 177 & \texttt{0R189} & \texttt{1R190} & 272 & \texttt{0R037} & \texttt{0L036} & 367 & \texttt{0R321} & \texttt{1R371} \\
83 & \texttt{0R042} & \texttt{1R097} & 178 & \texttt{0R191} & \texttt{1R192} & 273 & \texttt{0R037} & \texttt{1R284} & 368 & \texttt{0R081} & \texttt{1R372} \\
84 & \texttt{0L098} & \texttt{1R099} & 179 & \texttt{0R139} & \texttt{1R153} & 274 & \texttt{0R285} & \texttt{1R286} & 369 & \texttt{0R314} & \texttt{1R113} \\
85 & \texttt{0R097} & \texttt{1L100} & 180 & \texttt{0R193} & \texttt{1R194} & 275 & \texttt{0R189} & \texttt{1R249} & 370 & \texttt{0R217} & \texttt{1R373} \\
86 & \texttt{0R055} & \texttt{1R101} & 181 & \texttt{0R195} & \texttt{1R196} & 276 & \texttt{0R209} & \texttt{1R241} & 371 & \texttt{0R154} & \texttt{1R161} \\
87 & \texttt{0R081} & \texttt{1R102} & 182 & \texttt{0R050} & \texttt{1R197} & 277 & \texttt{0R287} & \texttt{0L036} & 372 & \texttt{0R374} & \texttt{0L036} \\
88 & \texttt{0R072} & \texttt{1R103} & 183 & \texttt{0R198} & \texttt{1R199} & 278 & \texttt{0R224} & \texttt{1R288} & 373 & \texttt{0R200} & \texttt{1R375} \\
89 & \texttt{0L104} & \texttt{1L038} & 184 & \texttt{0R200} & \texttt{1R201} & 279 & \texttt{0R072} & \texttt{1R289} & 374 & \texttt{0R143} & \texttt{1R376} \\
90 & \texttt{1R105} & \texttt{1R090} & 185 & \texttt{0R172} & \texttt{1R156} & 280 & \texttt{0R290} & \texttt{1R291} & 375 & \texttt{0R217} & \texttt{1R377} \\
91 & \texttt{0L106} & \texttt{1L091} & 186 & \texttt{0L202} & \texttt{1L140} & 281 & \texttt{0R292} & \texttt{1R293} & 376 & \texttt{0R209} & \texttt{1R328} \\
92 & \texttt{1R107} & \texttt{0R092} & 187 & \texttt{0R198} & \texttt{1R203} & 282 & \texttt{0R294} & \texttt{0L036} & 377 & \texttt{0R378} & \texttt{1R379} \\
93 & \texttt{0L108} & \texttt{1R078} & 188 & \texttt{0R187} & \texttt{1R204} & 283 & \texttt{0R295} & \texttt{0L036} & 378 & \texttt{0R169} & \texttt{1R128} \\
94 & \texttt{0L094} & \texttt{0L056} & 189 & \texttt{0L205} & \texttt{1L206} & 284 & \texttt{0R082} & \texttt{1L100} & 379 & \texttt{0R241} & \texttt{1R215} \\
\end{longtable}
\endgroup

\end{document}
